\section{Information changes the space of reasoning}

The previous chapter taught us how to construct probability models from
experiments. We learned to identify elementary outcomes, build the sample space
\(\Omega\), describe events as subsets of \(\Omega\), and assign probabilities
only after the model has been made explicit.

In this chapter we add a new idea: sometimes probability is not computed from
the original position of ignorance. Sometimes new information is given. When
that happens, the question is no longer asked inside the whole sample space.
The information restricts the set of possibilities that remain compatible with
what we know.

The central message is:
\[
\text{new information changes the space in which we reason.}
\]

This idea is the conceptual beginning of conditional probability. We shall not
begin by memorizing a formula. We shall first examine what information does to a
probability model.

\subsection*{Information removes impossible cases}

Suppose a die is rolled once. Before any information is given, the natural sample
space is
\[
\Omega=\{1,2,3,4,5,6\}.
\]
If the die is fair, all six elementary outcomes are equally likely.

Now imagine that someone tells us:
\[
\text{the result is even.}
\]
This information changes what is still possible. The outcomes
\[
1,3,5
\]
are no longer compatible with the information. The only outcomes that remain
possible are
\[
2,4,6.
\]
Thus we no longer reason inside the whole set \(\Omega\). We reason inside the
smaller event
\[
E=\{2,4,6\}.
\]

If we now ask
\[
\text{what is the probability that the result is greater than }3?
\]
we must not count favourable outcomes inside the original six outcomes. The
information has already told us that the result is even. Therefore the relevant
possibilities are only
\[
2,4,6.
\]
Among these, the outcomes greater than \(3\) are
\[
4,6.
\]
So, after receiving the information that the result is even, the probability of
``greater than \(3\)'' is obtained by comparing two favourable outcomes with the
three outcomes still possible.

\begin{remarkbox}
The old sample space has not disappeared from the model. It still explains what
the experiment was. But for the new question, after the information is known,
the active space of reasoning is the event compatible with that information.
\end{remarkbox}

\subsection*{A question may change its answer after information is given}

Before any information is given, the event
\[
A=\{4,5,6\}
\]
that the die result is greater than \(3\) has three favourable outcomes out of
six. In the fair die model, this gives probability \(1/2\).

After the information
\[
E=\{2,4,6\}
\]
is given, the same verbal question is evaluated in a smaller space. The relevant
outcomes are no longer
\[
1,2,3,4,5,6,
\]
but only
\[
2,4,6.
\]
Inside that restricted space, the event ``greater than \(3\)'' corresponds to
\[
A\cap E=\{4,6\}.
\]
Thus the answer changes. The event \(A\) itself did not change as a subset of
\(\Omega\). What changed is the set of outcomes that are still considered
possible.

This is the basic phenomenon of conditioning: information changes the reference
space.

\subsection*{Information as restriction to an event}

In probability, information usually has the form
\[
\text{we know that }B\text{ occurred.}
\]
The event \(B\) then becomes the new space of reasoning. Outcomes outside
\(B\) are no longer possible under the information.

If we are interested in another event \(A\), we do not ask how much of
\(\Omega\) belongs to \(A\). We ask how much of the still possible region \(B\)
also belongs to \(A\). The relevant part is therefore
\[
A\cap B.
\]

The geometry of the situation is:
\[
\Omega \quad\text{is the original universe,}
\]
\[
B \quad\text{is the information now known,}
\]
\[
A\cap B \quad\text{is the part of the information region where }A\text{ also holds.}
\]

\begin{definitionbox}
To \textbf{condition on information} means to restrict attention to the outcomes
compatible with that information and then measure the event of interest inside
that restricted region.
\end{definitionbox}

This is not yet a computational formula. It is the idea that the formula must
express.

\subsection*{Two dice: information about the sum}

Consider two fair dice, with order recorded. The original sample space is
\[
\Omega=\{(i,j):i,j\in\{1,2,3,4,5,6\}\}.
\]
It contains \(36\) ordered pairs.

Let
\[
B=\{\text{the sum is }7\}.
\]
Explicitly,
\[
B=\{(1,6),(2,5),(3,4),(4,3),(5,2),(6,1)\}.
\]
Suppose this information is given: the sum is \(7\). Then the six ordered pairs
above are the only outcomes still compatible with what we know.

Now ask:
\[
\text{what is the probability that the first die shows }5?
\]
Let
\[
A=\{\text{the first die shows }5\}.
\]
Inside the whole sample space, \(A\) contains six outcomes:
\[
(5,1),(5,2),(5,3),(5,4),(5,5),(5,6).
\]
But most of these are irrelevant once we know that the sum is \(7\). Inside the
restricted space \(B\), the only outcome with first coordinate \(5\) is
\[
(5,2).
\]
Thus, after the information is known, there is one favourable outcome among the
six outcomes still possible.

This example shows why one must be careful with the phrase ``the probability of
\(A\)''. Probability relative to what information? Before knowing the sum, the
question is evaluated in \(\Omega\). After knowing the sum, the question is
evaluated in \(B\).

\subsection*{The event is not replaced by the information}

A common misunderstanding is to think that once \(B\) is known, the event \(A\)
is replaced by \(B\). That is not correct. The event of interest remains \(A\).
The information is \(B\). The favourable outcomes are those where both things
are true:
\[
A\cap B.
\]

In the two-dice example:
\begin{itemize}
    \item \(B\) says that the sum is \(7\);
    \item \(A\) says that the first die is \(5\);
    \item \(A\cap B\) says that the first die is \(5\) and the sum is \(7\).
\end{itemize}

The information gives the new reference space. The intersection gives the part
of that reference space where the target event occurs.

\begin{remarkbox}
Conditioning always involves two roles: the event being asked about and the
information being conditioned on. Confusing these roles is one of the most
common sources of wrong answers.
\end{remarkbox}

\subsection*{Cards: the same card viewed after information}

Draw one card from a standard deck. Suppose the exact card is recorded, so the
sample space has \(52\) elementary outcomes.

Let
\[
A=\{\text{the card is an ace}\}.
\]
Before any additional information is given, there are \(4\) aces among \(52\)
cards.

Now suppose we are told:
\[
B=\{\text{the card is a heart}\}.
\]
The information reduces the relevant possibilities to the \(13\) hearts. Inside
this restricted set, only one card is an ace: the ace of hearts. Therefore, once
we know that the card is a heart, the event ``the card is an ace'' is evaluated
inside the set of hearts.

Now change the information. Suppose instead we are told:
\[
C=\{\text{the card is a face card}\}.
\]
The relevant possibilities are now the face cards. Inside that restricted set,
there are no aces, because aces are not face cards in the usual convention. Thus
the same event \(A\), ``the card is an ace'', behaves differently under different
information.

The event \(A\) did not change. The information changed the reference space.

\subsection*{Information may or may not change the probability}

New information restricts the space of reasoning, but it does not always change
the numerical probability of the event of interest.

For example, draw one card from a standard deck. Let
\[
A=\{\text{the card is an ace}\},
\]
and let
\[
B=\{\text{the card is a heart}\}.
\]
Knowing that the card is a heart leaves \(13\) possible cards, one of which is
an ace. The probability of being an ace inside the hearts is therefore the same
as the original proportion of aces in the whole deck:
\[
\text{one ace among thirteen hearts, just as four aces among fifty-two cards.}
\]
In this case the information ``heart'' does not change the probability of
``ace''.

But if the information is
\[
C=\{\text{the card is a face card}\},
\]
then the probability of being an ace becomes zero inside the restricted space.
This information does change the probability.

This observation prepares the idea of independence. Two events are independent
when learning that one occurred does not change the probability of the other.
But before defining independence formally, we must first understand what it means
to update the space of reasoning.

\subsection*{Renormalization: the whole information region becomes certain}

When the information \(B\) is given, the event \(B\) itself becomes certain in
the new situation. After being told that the die result is even, the statement
``the die result is even'' has probability \(1\) relative to the information.
After being told that the card is a heart, the statement ``the card is a heart''
has probability \(1\) relative to the information.

This means that probabilities inside the restricted region must be rescaled so
that the whole information region has total probability \(1\). In a uniform
finite model, this rescaling is easy to see. If \(B\) contains several equally
likely outcomes, then after we learn \(B\), those outcomes form the new complete
list of possibilities.

\begin{examplebox}
A fair die is rolled and we learn that the result is even. The restricted space
is
\[
B=\{2,4,6\}.
\]
Inside this new space, the three outcomes \(2,4,6\) are treated as the complete
possibilities. Therefore the event \(B\) itself has probability \(1\) relative to
this information, and each of the three remaining outcomes has probability
\(1/3\) in the restricted uniform model.
\end{examplebox}

In non-uniform models, the same idea holds, but the remaining probabilities must
be rescaled proportionally. The original weights outside \(B\) are discarded
because they correspond to outcomes no longer compatible with the information.
The weights inside \(B\) are then adjusted so that their total becomes \(1\).

This rescaling principle is the reason the formal definition of conditional
probability has the structure it has. The numerator must measure the part of the
information region where the target event occurs. The denominator must measure
the whole information region, because that whole region becomes the new unit of
probability.

\subsection*{Why the information must have positive probability}

Conditioning on information only makes sense if the information can occur in the
model. If \(P(B)=0\), then the instruction ``reason inside \(B\)'' cannot be
implemented by rescaling probabilities inside \(B\), because there is no positive
probability mass to rescale.

In finite elementary examples this is easy to understand. Suppose a standard die
is rolled and someone says:
\[
\text{the result is }8.
\]
This statement corresponds to the empty event in the die model
\(\Omega=\{1,2,3,4,5,6\}\). There are no outcomes compatible with it. It cannot
become a new space of reasoning.

More generally, conditional probability is defined only when the conditioning
event has positive probability. This condition is not a technical decoration. It
says that the information must leave a meaningful region in which probability can
be redistributed.

\subsection*{The conceptual picture}

The logic of conditioning can now be summarized as follows.

\begin{enumerate}
    \item We begin with a probability model on \(\Omega\).
    \item We receive information that an event \(B\) occurred.
    \item Outcomes outside \(B\) are no longer compatible with the information.
    \item The event \(B\) becomes the new reference space.
    \item To study another event \(A\), we look at the part of \(B\) where \(A\)
    is also true, namely \(A\cap B\).
    \item Probabilities inside \(B\) are rescaled so that \(B\) has total
    probability \(1\) in the new situation.
\end{enumerate}

\begin{theorembox}
The guiding principle of conditional probability is:
\[
\text{probability after information is probability inside the event described by that information.}
\]
The formal definition in the next section will turn this principle into a
precise rule.
\end{theorembox}

The important point is that the formula will not be arbitrary. It must express
the geometry we have already seen: restrict to the information event, keep only
the part where the target event also occurs, and renormalize the information
region so that it becomes the new whole space.
